\Needspace{11\baselineskip}
\SourceLeaf{277}
\begin{center}
\textit{Composition of classes.}
\end{center}

\Article{249}
\textsc{Theorem.} \textit{If the forms $f,f'$ belong to the same orders, genera and classes as $g,g'$ respectively, the form composed from $f$ and $f'$ belongs to the same class as the form composed from $g$ and $g'$.}

This theorem, whose truth follows at once from Article~239, gives the meaning of a \textit{class composed from two given classes}, or even \textit{from several}.

If any class $K$ is composed with the principal class, the class $K$ itself results; that is, the principal class can be left out when composing it with other classes of the same determinant. Composing two opposite properly primitive classes always gives the principal class of the same determinant (see Article~243). Since every ambiguous class is opposite to itself, composing any properly primitive ambiguous class with itself gives the principal class of the same determinant.

The converse of the last proposition also holds: namely, \textit{if composing a properly primitive class $K$ with itself gives the principal class $H$ of the same determinant, $K$ must be an ambiguous class.} For if $K'$ is the class opposite to $K$, the three classes $K,K,K'$ give the same composite class as $H$ and $K'$. The former give $K$ (since $K$ and $K'$ produce $H$, and this with $K$ gives $K$ itself), the latter $K'$; so $K$ coincides with $K'$ and is therefore an ambiguous class.

The following proposition should also be noted: \textit{If the classes $K,L$ are opposite to the classes $K',L'$ respectively, the class composed from $K$ and $L$ is opposite to the class composed from $K'$ and $L'$.} Let the forms $f,g,f',g'$ belong to the classes $K,L,K',L'$ respectively; let the form $F$ be composed from $f,g$, and $F'$ from $f',g'$. Since $f'$ is improperly equivalent to $f$, and $g'$ to $g$, while $F$ is composed directly from each of $f,g$, $F$ is also composed from $f',g'$, but inversely from each. So any form improperly equivalent to $F$ is composed directly from $f',g'$ and is therefore properly equivalent to $F'$ (Articles~238,~239). Thus $F,F'$ are improperly equivalent, and the classes to which they belong are opposite.

It follows that if an ambiguous class $K$ is composed with an ambiguous class $L$. an ambiguous class always results. For it is opposite to the class composed from the classes opposite to $K,L$, and therefore to itself, since these classes are opposite to themselves.

\SourceLeaf{278}
Finally, we observe that, given any two classes $K,L$ of the same determinant, the first of which is properly primitive, we can always find a class $M$ of the same determinant whose composition with $K$ gives $L$. Clearly we achieve this by taking for $M$ the class composed from $L$ and the class opposite to $K$. It is also very easy to see that this is the only class with this property; that is, distinct classes of the same determinant, composed with the same properly primitive class, produce distinct classes.

Composition of classes can conveniently be denoted by the addition sign $+$, just as identity of classes is denoted by the equality sign. With these signs the proposition just given can be expressed as follows. If $K'$ is the class opposite to $K$, $K+K'$ is the principal class of the same determinant, so $K+K'+L=L$; setting $K'+L=M$ therefore gives $K+M=L$, as required. But if, besides $M$, there were another $M$ with the same property, that is, $K+M'=L$, we would have $K+K'+M'=L+K'=M$, giving $M'=M$. --- If several identical classes are composed, we can indicate this, as in multiplication, by putting their number before them, so that $2K$ means the same as $K+K$, $3K$ the same as $K+K+K$ etc. The same signs could also be applied to forms, so that $(a,b,c)+(a',b',c')$ would denote a form composed from $(a,b,c)$, $(a',b',c')$. But to avoid even the appearance of ambiguity, we prefer not to use this abbreviation, especially since we have already assigned a particular meaning to a sign such as $\surd M(a,b,c)$. --- We say that the class $2K$ arises by \textit{doubling} the class $K$, the class $3K$ by \textit{tripling}, and so on.

\Needspace{8\baselineskip}
\Article{250}
If $D$ is a number divisible by $mm$ (where we suppose $m$ positive), there is an order of forms of determinant $D$ derived from the properly primitive order of determinant $\frac{D}{mm}$ (or \textit{two} when $D$ is negative, namely a positive and a negative order). Clearly the form $(m,0,-\frac{D}{m})$ belongs to that order, namely the positive one, and may rightly be regarded as its \textit{simplest form} (just as $(-m,0,\frac{D}{m})$ is the simplest in the negative order when $D$ is negative). If also $\frac{D}{mm}\equiv1$ (modulo $4$), there is an order of forms of determinant $D$ derived from the improperly primitive order of determinant $\frac{D}{mm}$. Clearly the form $(2m,m,\frac{mm-D}{2m})$ belongs to it and is regarded as its simplest form. (When $D$ is negative there are again two orders, and in the negative order the form $(-2m,-m,\frac{D-mm}{2m})$ is regarded as the simplest.) Thus, \textit{for example,} if we also include the case where $m=1$, in the four orders
\SourceLeaf{279}
of forms of determinant $45$ the simplest forms are $(1,0,-45)$, $(2,1,-22)$. $(3,0,-15)$, $(6,3,-6)$. With this understood, we have the following

\textsc{Problem.} \textit{Given any form $F$ of the order $O$, find a properly primitive (positive) form of the same determinant whose composition with the simplest form in $O$ gives $F$.}

\textit{Solution.} Let the form $F=(ma,mb,mc)$ be derived from the primitive form $f=(a,b,c)$, whose determinant is $=d$, and suppose \textit{first} that $f$ is properly primitive. First we observe that if $a$ is not coprime to $2dm$, there are certainly other forms properly equivalent to $(a,b,c)$ whose leading coefficients have this property. For by Article~228 that form represents numbers coprime to $2dm$. Let $a'=a\alpha\alpha+2b\alpha\gamma+c\gamma\gamma$ be such a number, and suppose, as we may, that $\alpha,\gamma$ are coprime. Then choose $\beta,\delta$ so that $\alpha\delta-\beta\gamma=1$, and transform $f$ by the substitution $\alpha,\beta,\gamma,\delta$ into the form $(a',b',c')$, which is properly equivalent to it and has the required property. Since $F$ and $(a'm,b'm,c'm)$ are also properly equivalent, we easily see that it is enough to consider the case where $a$ is coprime to $2dm$. Then $(a,bm,cmm)$ is a properly primitive form (for if $a,2bm,cmm$ had a common divisor, it would also divide $2dm=2bbm-2acm$) of the same determinant as $F$. It is easily confirmed that $F$ is transformed into the product of the form $(m,0,-dm)$, which, unless $F$ is a negative form, is the simplest in the order $O$, and $(a,bm,cmm)$, by the substitution $1,0,-b,-cm;\ 0,m,a,bm$. So the criterion in observation~4 of Article~235 shows that $F$ is composed from $(m,0,-dm)$ and $(a,bm,cmm)$. But when $F$ is a negative form, it is transformed into the product of the simplest form of the same order $(-m,0,dm)$ and the positive form $(-a,bm,-cmm)$ by the substitution $1,0,b,-cm;\ 0,-m,-a,bm$, and is therefore composed from them.

\textit{Second}, if $f$ is an improperly primitive form, we may suppose $\frac{1}{2}a$ coprime to $2dm$. For if this property does not yet hold for the form $f$, a form properly equivalent to $f$ and having this property can be found. It follows easily that $(\frac{1}{2}a,bm,2cmm)$ is a properly primitive form of the same determinant as $F$. We confirm just as easily that $F$ is transformed into the product of the forms
\[
(\pm2m,\pm m,\pm\tfrac{1}{2}(m-dm)),\quad
(\pm\tfrac{1}{2}a,bm,\pm2cmm)
\]
by the substitution
\SourceLeaf{280}
\[
1,0,\tfrac{1}{2}(1\mp b),-cm;\quad 0,\pm2m,\pm\tfrac{1}{2}a,(b\pm1)m
\]
where the lower signs are to be taken when $F$ is a negative form and the upper signs in the remaining cases. The form is therefore composed from these two forms, the former being the simplest in the order $O$, the latter a properly primitive (positive) form.

\Needspace{8\baselineskip}
\Article{251}
\textsc{Problem.} \textit{Given two forms $F,f$ of the same determinant $D$ and belonging to the same order $O$, find a properly primitive form of determinant $D$ whose composition with $f$ produces $F$.}

\textit{Solution.} Let $\varphi$ be the simplest form of the order $O$; let $\mathfrak{F}$, $\mathfrak{f}$ be properly primitive forms of determinant $D$ whose compositions with $\varphi$ give $F.\ f$ respectively; finally, let $f'$ be a properly primitive form whose composition with $\mathfrak{f}$ gives $\mathfrak{F}$. Then the form $F$ is composed from the three forms $\varphi,\mathfrak{f},f'$, or from the two $f,f'$. \textit{Q.~E.~I.}

Thus any class of a given order can be regarded as composed from any given class of the same order and some properly primitive class of the same determinant.

\Needspace{11\baselineskip}
\begin{center}
\textit{For a given determinant, the individual genera of the same order contain equally many classes.}
\end{center}

\Article{252}
\textsc{Theorem.} \textit{For a given determinant, the individual genera of the same order contain equally many classes.}

\textit{Proof.} Let the genera $G$ and $H$ belong to the same order, let $G$ consist of $n$ classes $K,K',K''\ldots K^{n-1}$, and let $L$ be any class of the genus $H$. By the preceding article, find a properly primitive class $M$ of the same determinant whose composition with $K$ gives $L$, and denote the classes obtained by composing $M$ with $K',K''\ldots K^{n-1}$ by $L',L''\ldots L^{n-1}$ respectively. Then the last observation of Article~249 shows that all the classes $L,L',L''\ldots L^{n-1}$ are distinct, and by Article~248 they all belong to the same genus, \textit{that is,}\ the genus $H$. Finally, we easily see that $H$ cannot contain any other classes, since any class of the genus $H$ can be regarded as composed from $M$ and another class of the same determinant, which must always belong to the genus $G$. So $H$, like $G$, contains $n$ distinct classes. \textit{Q.~E.~D.}
